\documentclass[11pt]{article}
\usepackage{../homework}

\classcode{MATH 465 Differential Geometry}
\profname{Dr. David Singer}
\duedate{16 September 2026}
\assignnum{Assignment 03}

\begin{document}
All problems from Do Carmo, \textit{Differential Geometry of Curves and Surfaces} 2e.

We recall the statements of some propositions here for later use.

\begin{prop}
	\label{prop:1}
	If \( f: U \rightarrow \mathbb{R} \) is a differentiable function in an open set \( U \) of \( \mathbb{R}^2 \), then the graph of \( f \), that is, the subset of \( \mathbb{R}^3 \) given by \( (x,y, f(x,y)) \) for \( (x,y) \in U \), is a regular surface.
\end{prop}

\begin{prop}
	\label{prop:2}
	If \( f: U \subset \mathbb{R}^3 \rightarrow \mathbb{R} \) is a differentiable function and \( a \in f(U) \) is a regular value of \( f \), then \( f^{-1}(a) \) is a regular surface in \( \mathbb{R}^3 \).
\end{prop}

\begin{problem}{2.2.6}
	Give another proof of Proposition 1 by applying Proposition 2 to \( h(x,y,z) = f(x,y) - z \).
\end{problem}
\begin{proof}
	Let \( h: U \subset \mathbb{R}^3 \rightarrow \mathbb{R} \) be given by \( h(x,y,z) = f(x,y) - z \), for some differentiable \( f: V \subset \mathbb{R}^2 \rightarrow \mathbb{R} \) defined on an open set \( V \).
	Clearly \( h \) is differentiable, with differential
	\begin{equation}
		\label{eq:6}
		dh_{p} = \begin{pmatrix}
			\partial_xf(x,y) & \partial_yf(x,y) & -1
		\end{pmatrix}.
	\end{equation}

	This is clearly never singular.
	In particular, \( 0 \) is a regular value, so the preimage \( h^{-1}(0) = \{(x,y,z) \in \mathbb{R}^3 \mid f(x,y) - z = 0\}\) is a regular surface in \( \mathbb{R}^3 \) by \cref{prop:2}.
	Now, for any \( (x,y) \in V \), the point \( (x,y, f(x,y)) \in h^{-1}(0) \) by construction, and this proves \cref{prop:1}.
\end{proof}

\begin{problem}{2.2.12}
	Show that \( \mathbf{x}: U \subset \mathbb{R}^2 \rightarrow \mathbb{R}^3 \) given by
	\begin{equation}
		\label{eq:1}
		\mathbf{x}(u,v) = (a\sin u\cos v, b \sin u \sin v, c \cos u), \quad a,b,c \neq 0,
	\end{equation}
	where \( 0 < u < \pi, 0 < v < 2\pi \), is a parametrization for the ellipsoid
	\begin{equation}
		\label{eq:2}
		\frac{x^2}{a^2} + \frac{y^2}{b^2} + \frac{z^2}{c^2} = 1.
	\end{equation}
\end{problem}
\begin{proof}
	To begin with, \( \mathbf{x} \) is obviously differentiable (it has continuous partial derivatives of all orders), and in fact
	\begin{equation}
		\label{eq:7}
		d \mathbf{x}_p = \begin{pmatrix}
			a\cos  u \cos  v  & -a \sin  u \sin  v \\
			b \cos  u \sin  v & b \sin  u \cos  v  \\
			-c \sin  u        & 0
		\end{pmatrix}.
	\end{equation}

  To show that \( d \mathbf{x}_p \) is injective it suffices to show the columns are linearly independent, that is, that they have non-zero cross product.
  \begin{align}
  \label{eq:10}
  d \mathbf{x}_p(e_1) \times d \mathbf{x}_p(e_2) &= \det \begin{pmatrix}
\hat{x}             & \hat{y} & \hat{z} \\
a \cos  u \cos  v   & b \cos  u \sin  v & -c \sin  u \\
-a \sin  u \sin  v  & b \sin  u \cos  v & 0
  \end{pmatrix} \\
  &= \left( bc \sin^2u \cos v, ac \sin^2 u \sin v, ab \cos u \sin u\right)^{T}.
  \end{align}
  This would only vanish when \( \sin u = 0 \) which cannot happen in our domain of \( U = (0, \pi) \times (0, 2\pi) \).

  Finally, we show that \( \mathbf{x}(U) \subset S \) where \( S \) is the ellipsoid, and that \( \mathbf{x} \) is a homeomorphism onto its image.
  Let \( (u,v) \in U \), so that \( \mathbf{x}(u,v) = (a\sin u \cos v, b \sin u \sin v, c \cos u) \).
  Then
  \begin{equation}
  \label{eq:11}
   \frac{x^2}{a^2} + \frac{y^2}{b^2} + \frac{z^2}{c^2} = \sin^2 u \cos^2 v + \sin^2 u \sin^2 v + \cos^2 u = \sin^2 u + \cos^2 u = 1
  \end{equation}
  and thus \( \mathbf{x}(U) \in S \).

  \( z = c\cos u \) uniquely determines \( u \in (0, \pi)\), while, for fixed \( u \), \( x= a \sin u \cos v \) and \( y = a\cos u \cos v \) uniquely determines \( v \in (0, 2\pi) \), so \( \mathbf{x} \) is injective on its image and is therefore a homeomorphism.
\end{proof}

\begin{problem}{2.2.16}
	Consider the stereographic projection \(\pi: S^2 \setminus \{(0,0,2)\} \rightarrow \mathbb{R}^2\) of the unit sphere centered at \( (0,0,1) \), that is, the set given by
	\begin{equation}
		\label{eq:3}
		x^2 + y^2 + (z-1)^2 = 1,
	\end{equation}
	which carries a point \( p = (x,y,z) \) of the sphere without the north pole onto the intersection of the \( xy \) plane with the straight line with connects the north pole to \( p \).
	Let \( (u,v) = \pi(x,y,z) \).
	\begin{enumerate}[label=(\alph*)]
		\item Show that \( \pi^{-1}: \mathbb{R}^2 \rightarrow S^2 \) is given by
		      \begin{equation}
			      \label{eq:4}
			      \pi^{-1}(u,v) = \left( \frac{4u}{u^2 v^2 + 4}, \frac{4v}{u^2 + v^2 + 4}, \frac{2(u^2 + v^2)}{u^2 + v^2 + 4} \right).
		      \end{equation}
          \begin{proof}
            \(\pi^{-1}\) maps a point \( (u,v) \in \mathbb{R}^2 \) to a point \( (x,y,z) \in S^2 \setminus \{\text{north pole}\} \).
            Given a point \( (u,v) \), the line from to the north pole can be parametrized as
            \begin{equation}
            \label{eq:12}
            \alpha(t) = (1-t)(0,0,2) + t(u, v, 0) = (tu, tv, 2 - 2t).
            \end{equation}
            We then ask when \(\alpha(t)\) intersects the sphere, that is,
            \begin{align}
            \label{eq:13}
              &(tu)^2 + (tv)^2 + (2 - 2t - 1)^2 = 1 \\
              &\Leftrightarrow t^2(u^2 + v^2) + 4t^2 - 4t = 0 \\
              &\Leftrightarrow t^2(u^2 + v^2 + 4) = 4t \\
              &\Leftrightarrow t = \frac{4}{u^2 + v^2 + 4}.
            \end{align}
            Substituting this value for \( t \) into \cref{eq:12} gives the point at which the line from the north pole to the point \( (u,v) \) intersects the sphere, that is,
            \begin{align}
            \label{eq:14}
              \pi^{-1}(u,v) &= \left( \frac{4u}{u^2 + v^2 + 4}, \frac{4v}{u^2 + v^2 + 4}, 2(1 - \frac{4}{u^2 + v^2 + 4} ) \right) \\
              & = \left( \frac{4u}{u^2 + v^2 + 4}, \frac{4v}{u^2 + v^2 + 4}, \frac{2(u^2 + v^2)}{u^2 + v^2 + 4} \right),
            \end{align}
            as desired.
          \end{proof}
		\item Show that it is possible, using stereographic projection, to cover the sphere with two coordinate neighborhoods.
          \begin{proof}
            Our two coordinate neighborhoods will be \( S^2 \setminus \{\text{north pole}\} \) and \( S^2 \setminus \{\text{south pole}\} \), where for the latter we define a symmetric stereographic projection from the south pole.
            We then must check that the overlap, \( S^2 \setminus \{\text{north and south poles}\} \) has diffeomorphic transition maps.

            Let the south pole be \( S = (0,0,0) \).
            We can define a second coordinate neighborhood \( S^2 \setminus \{S\} \) using a stereographic projection \( \pi_2: S^2 \setminus \{S\} \rightarrow \mathbb{R}^2 \) that projects a point \( p = (x,y,z) \) onto the plane \( z = 2 \).

            The straight line connecting \( S \) to \( p \) is parametrized by \( \beta(s) = s(x,y,z) = (sx, sy, sz) \).
            This line intersects the plane \( z = 2 \) when \( sz = 2 \), which implies \( s = \frac{2}{z} \).
            Substituting this back gives the projection:
            \begin{equation}
                \pi_2(x,y,z) = \left( \frac{2x}{z}, \frac{2y}{z} \right).
            \end{equation}

            The overlap between the two neighborhoods is the sphere minus both poles, \( W = S^2 \setminus \{(0,0,2), (0,0,0)\} \).
            Under the first projection \( \pi_1 \), this domain corresponds to \( \mathbb{R}^2 \setminus \{(0,0)\} \).
            To verify compatibility, we must show that the transition map \( \pi_2 \circ \pi_1^{-1}: \mathbb{R}^2 \setminus \{(0,0)\} \rightarrow \mathbb{R}^2 \setminus \{(0,0)\} \) is a diffeomorphism.

            Using the formula for \( \pi_1^{-1}(u,v) \) from part (a), we substitute its \( x \), \( y \), and \( z \) coordinates into \( \pi_2 \).
            \begin{align}
                \pi_2 \circ \pi_1^{-1}(u,v) &= \pi_2\left( \frac{4u}{u^2 + v^2 + 4}, \frac{4v}{u^2 + v^2 + 4}, \frac{2(u^2 + v^2)}{u^2 + v^2 + 4} \right) \\
                &= \left( \frac{2 \left( \frac{4u}{u^2 + v^2 + 4} \right)}{\frac{2(u^2 + v^2)}{u^2 + v^2 + 4}}, \frac{2 \left( \frac{4v}{u^2 + v^2 + 4} \right)}{\frac{2(u^2 + v^2)}{u^2 + v^2 + 4}} \right) \\
                &= \left( \frac{4u}{u^2 + v^2}, \frac{4v}{u^2 + v^2} \right).
            \end{align}

            Since \( u^2 + v^2 \neq 0 \) for all \( (u,v) \in \mathbb{R}^2 \setminus \{(0,0)\} \), the components of this transition map are rational functions with non-vanishing denominators.
            Therefore, \( \pi_2 \circ \pi_1^{-1} \) is smooth (\( C^\infty \)).

            Because of the geometric symmetry of the sphere and the chosen projection planes, the inverse transition map \( \pi_1 \circ \pi_2^{-1} \) takes the exact same form, and is also smooth on \( \mathbb{R}^2 \setminus \{(0,0)\} \).
            Since the transition map and its inverse are both smooth, they are diffeomorphisms, and our claim follows from here.
          \end{proof}
	\end{enumerate}
\end{problem}

\begin{problem}{2.3.3}
	Show that the paraboloid \( z = x^2 + y^2 \) is diffeomorphic to a plane.
\end{problem}
\begin{proof}
Let \( \Sigma \) be the surface of the paraboloid.
We want to show there exists a differential map \( \varphi: \mathbb{R}^2 \rightarrow \Sigma \) with differentiable inverse \( \varphi^{-1}: \Sigma \rightarrow \mathbb{R}^2 \).

Consider the natural parametrization
\begin{equation}
\label{eq:15}
\varphi(x,y) = (x, y, x^2 + y^2).
\end{equation}

This is surjective, by construction.
It is also injective, since if \( (x_1, y_1, x_1^2 + y_1^2) = (x_2, y_2, x_2^2 + y_2^2) \), it is obviously true that \( (x_1, y_1) = (x_2, y_2) \).
Moreover, it is clearly differentiable.

Let us claim that \( \varphi^{-1}(x, y, x^2 + z^2) = (x, y) \) is a differentiable inverse; bijectivity and differentiability follow immediately and identically to above.
This is the inverse, since
\((\varphi \circ \varphi^{-1})(x,y, x^2 + y^2) = \varphi(x,y) = (x, y, x^2 + y^2) = \text{Id}(x, y, x^2 + y^2) \), and \( (\varphi^{-1} \circ \varphi)(x,y) = \varphi^{-1}(x, y, x^2 + y^2) = (x, y) = \text{Id}(x, y) \).
\end{proof}

\begin{problem}{2.3.12}
	Parametrized surfaces are often useful to describe sets \(\Sigma\) which are regular except for a finite number of points and a finite number of lines.
	For instance, let \( C \) be the trace of a regular parametrized curve \(\alpha: (a,b) \rightarrow \mathbb{R}^3\) which does not pass through the origin \( O = (0,0,0) \).
	Let \(\Sigma\) be the set generated by the displacement of a straight line \( l \) passing through a moving point \( p \in C \) and the fixed point \( 0 \) (a cone with vertex \( 0 \)).
	\begin{enumerate}[label=(\alph*)]
		\item Find a parametrized surface \( \mathbf{x} \) whose trace is \(\Sigma\).
          \begin{proof}
            Writing the parameter for the curve \(\alpha\) as \( s \in (a,b) \), points in the trace \(\Sigma\) are of the form \( t\alpha(s) + (1-t)(0,0,0) = t \alpha(s) \).
            Hence, if we define
            \begin{equation}
            \label{eq:16}
            \mathbf{x}: (a,b) \times \mathbb{R} \subset \mathbb{R}^2 \rightarrow \mathbb{R}^3, \quad (s, t) \mapsto t \alpha(s),
            \end{equation}
            then \( \mathbf{x} \) is clearly differentiable as a scalar multiple of a differentiable function, and \( \mathbf{x}(U) = \Sigma \) where \( U = (a,b) \times \mathbb{R} \).
          \end{proof}
		\item Find the points where \( \mathbf{x} \) is not regular.
          \begin{proof}
            The differential is the \( 3 \times 2 \) matrix whose columns are the partial derivatives \( \partial_s\mathbf{x}, \partial_t\mathbf{x} \), and thus \( \mathbf{x} \) is regular at \( p \) if and only if \(\partial_s\mathbf{x} \times \partial_t\mathbf{x} = t(\alpha^{\prime}(s) \times \alpha(s))\neq 0\).

            Note that this fails when \( t=0 \), that is, for all points \( p \in \mathbf{x}^{-1}(0) \) which map to the origin.
            The other singularities occur when
            \begin{align}
            \label{eq:8}
              &\alpha(s) \times \alpha^{\prime}(s) = 0 \\
              \Rightarrow \ &\alpha^{\prime}(s) = \lambda\alpha(s),
            \end{align}
            for some \( \lambda \in \mathbb{R} \).
          \end{proof}
		\item What should be removed from \(\Sigma\) so the remaining set is a regular surface?
          \begin{proof}
            We must first of all remove the singular points from part (b), that is,
            \begin{equation}
             \{ \mathbf{x}(s,t) \mid t = 0 \text{ or } \alpha^{\prime}(s) = \lambda\alpha(s), \lambda \in \mathbb{R} \}.
            \end{equation}
            However, since we are only covering the surface with a single coordinate patch, we need \( \mathbf{x} \) to be a global homeomorphism, or, in particular, injective.
            Thus, we must also remove any points of self intersection,
            \begin{equation}
            \label{eq:9}
            \{\mathbf{x}(s,t) \mid \mathbf{x}(s,t) = \mathbf{x}(s^{\prime}, t^{\prime}), (s, t) \neq (s^{\prime}, t^{\prime}) \in \mathbb{R}^2\}.
            \end{equation}
          \end{proof}
	\end{enumerate}
\end{problem}
\end{document}
